\documentclass[a4paper,10pt]{article}
\usepackage[utf8]{inputenc}
\usepackage[margin=1.5cm, bottom=2cm]{geometry}
\usepackage{fancyhdr}
\usepackage{enumitem}
\usepackage{csquotes}
\usepackage[utf8]{inputenc}
\usepackage[italian]{babel}
\usepackage[T1]{fontenc}

\usepackage{tikz}
\usetikzlibrary{automata, positioning}
\pagestyle{fancy}
\setlength{\parindent}{0pt}
\usepackage[most]{tcolorbox}

\newtcolorbox{osservazioni}{
  colback=gray!8,    % sfondo grigio chiaro
  colframe=white,      % nessun bordo
  arc=2mm,            % angoli smussati
  boxrule=0pt,        % spessore bordo 0
  fontupper=\small,
  left=5mm, right=5mm, top=3mm, bottom=3mm,
  before skip=5pt,   % niente spazio prima
  after skip=5pt     % niente spazio dopo
}
\pagestyle{fancy}


%%%%%%%%%%%%%%%%%%%
\newcommand{\EsercitazioneNum}{01}
\newcommand{\Data}{03/10/2025}
%%%%%%%%%%%%%%%%%%%

% Numero di pagina in basso al centro
\fancyhf{}
\fancyfoot[C] \thepage
\renewcommand{\footrulewidth}{0pt} % niente linea in basso
\renewcommand{\headrulewidth}{0pt} % niente linea in alto

\begin{document}
\footnotesize  % font generale più piccolo

\noindent
\small \textbf{Computabilità, Complessità e Logica - a.a. 2025/26} \hfill \textbf{Approfondimento \EsercitazioneNum} \\[0.2cm]
\scriptsize Matteo Liotta, \texttt{matteo.liotta@studenti.units.it} \hfill \Data \\[0.2cm]

\noindent\makebox[\linewidth]{\rule{\linewidth}{0.4pt}} \\[1cm]


\noindent\Large \textbf{Approfondimento \EsercitazioneNum}: \textit{Non determinismo e "scommesse"}\\[0cm]
%\\[0.4cm]

\small

\underline{\textbf{Introduzione al problema}}\\[0.3cm]
Introducendo gli automi non deterministici, si è proposto a lezione un esempio tale da mostrare il vantaggio derivato dall'impiego del non determinismo, in termini di complessità di scrittura dell'automa che riconosca un linguaggio desiderato.

Si può considerare dunque il seguente lunguaggio su $\Sigma=\{a,b,c\}$, in cui è evidente tale comportamento:

$$L_n=\{w\#w'| w,w'\in\Sigma^*, \quad \exists x\in \Sigma \quad t.c.\quad  x\in w'\ e\ x\in w\}$$\\[0.01cm]
Questo linguaggio è tale da contenere qualsiasi parola della forma $$w\#w'$$ in cui ambo le parti $w$ e $w'$ contengano \textbf{almeno un simbolo} $x\in \Sigma$ in comune. Consideriamo per semplicità il caso in cui $$|\Sigma|=n=3$$

\underline{\textbf{Caso deterministico}}\\[0.3cm]
Come visto durante una esercitazione (cfr. 01), ogni stato di un automa regolare può essere utilizzato in modo da rappresentare conoscenza intorno alla situazione corrente; in questo caso, possiamo utilizzare ogni stato in modo da evidenziare se alcune lettere sono state incontrate all'interno di una parola: 

\begin{itemize}
    \item $aabbaba\# \rightarrow \{a,b\} $, in quanto vengono letti entrambi i simboli $a$ e $b$, almeno una volta
    \item $a \rightarrow \{a\}$
    \item $ab \rightarrow \{a,b\}$
\end{itemize}
È dunque possibile iniziare con $n$ transizioni dallo stato iniziale, una per simbolo di $\Sigma$, così da raggiungere gli stati che rappresentano l'insieme dei simboli letti almeno una volta, della forma $\{a\}, \{b\} $ o $\{c\}$.

\begin{osservazioni}
    \textit{\textbf{Nota}}. \textit{Si noti come non è richiesta alcuna capacità di "contare" all'automa, bensì memorizzare se è avvenuta \underline{almeno una} occorrenza di un simbolo all'interno di una parola}. Quando considereremo la lettura di $w'$, sbasterà osservare se almeno una delle lettere presenti -almeno una volta- in $w$ vengono trovate.
\end{osservazioni}
A questo punto, dato $$\Sigma=\{a,b,c\}$$ possiamo considerare, da uno stato iniziale, \textbf{tre percorsi}, uno relativo alla lettura di ciascuna lettera di $\Sigma$. 

{
\small
\begin{center}
\begin{tikzpicture}[shorten >=1pt, node distance=2cm, on grid, auto]
    \node[state, initial, initial text=] (init)   {$\emptyset$};
    \node[state, right=of init] (b) {$\{b\}$};
    \node[state, above=1.5cm of b] (a) {$\{a\}$};
    \node[state, below=1.5cm of b] (c) {$\{c\}$};

    \path[->]
    (init) edge node {a} (a)
    (init) edge node {b} (b)
    (init) edge node {c} (c);
    
\end{tikzpicture}
\end{center}
}

Ogni stato raggiunto memorizza l'informazione relativa alla lettura del corrispondente simbolo. Poiché di questi non è necessario memorizzare ulteriori occorrenze, possiamo introdurre dei loop per garantire generalità:

\vfill
    \noindent\makebox[\linewidth]{\rule{\linewidth}{0.4pt}}
    \newpage
    {\noindent
    \small \textbf{Computabilità, Complessità e Logica - a.a. 2025/26} \hfill \textbf{Approfondimento \EsercitazioneNum} \\[0.2cm]
    \scriptsize Matteo Liotta, \texttt{matteo.liotta@studenti.units.it} \hfill \Data \\[0.2cm]
    \noindent\makebox[\linewidth]{\rule{\linewidth}{0.4pt}} \\[0.3cm]}

{
\small
\begin{center}
\begin{tikzpicture}[shorten >=1pt, node distance=2cm, on grid, auto]
    \node[state, initial, initial text=] (init)   {$\emptyset$};
    \node[state, right=of init] (b) {$\{b\}$};
    \node[state, above=1.5cm of b] (a) {$\{a\}$};
    \node[state, below=1.5cm of b] (c) {$\{c\}$};

    \path[->]
    (init) edge node {a} (a)
    (init) edge node {b} (b)
    (init) edge node {c} (c)
    (a) edge [loop left, looseness=3] node {a} (a)
    (c) edge [loop left, looseness=3] node {c} (c)
    (b) edge[loop above, looseness=3] node {b} (b);
    
\end{tikzpicture}
\end{center}
}

In ciascun caso, ora, possiamo procedere e aspettarci eventualmente la lettura di altri termini di $\Sigma$ non ancora incontrati, come $b$ e $c$ nel caso di avvenuta lettura di sole $a$. Questo può, in linea generale, portarci ad aspettarci  stati quali tutte le combinazioni possibili di elementi di $\Sigma$:$$\{a\}, \{b\}, \{c\}, \{a,b\}, \{b,c\}, \{a,c\}, \{a,b,c\}$$
Quindi, l'automa può essere aggiornato introducendo:

{
\small
\begin{center}
\begin{tikzpicture}[shorten >=1pt, node distance=2cm, on grid, auto]
    \node[state, initial, initial text=] (init) {$\emptyset$};
    \node[state, right=of init] (b) {$\{b\}$};
    \node[state, above=of b] (a) {$\{a\}$};
    \node[state, below=of b] (c) {$\{c\}$};
    
    \node[state, right=of a] (ab) {$\{a,b\}$};
    \node[state, right=of b] (bc) {$\{b,c\}$};
    \node[state, right=of c] (ac) {$\{a,c\}$};

    \node[state, right=of bc] (abc) {$\{a,b,c\}$};
    
    \path[->]
    (init) edge node {a} (a)
    (init) edge node {b} (b)
    (init) edge node {c} (c)
    (a) edge [loop left, looseness=3] node {a} (a)
        edge node {b} (ab)
        edge [bend left=3] node {c} (ac)
    (c) edge [loop left, looseness=3] node {c} (c)
        edge node {b} (bc)
        edge [bend right] node {a} (ac)
    (b) edge[loop above, looseness=3] node {b} (b)
        edge [bend right] node {a} (ab)
        edge [bend right] node {c} (bc)
    

    (ab) edge [loop above, looseness=3] node {a,b} (ab)
         edge node {c} (abc)
    (bc) edge [loop above, looseness=3] node {b,c} (bc)
    edge node {a} (abc)
    (ac) edge[loop above, looseness=3] node {a,c} (ac)
    edge node {b} (abc);
    
\end{tikzpicture}
\end{center}
}

La prima parte è completata: osservata la componente $w$ di $w\#w'$ ci aspettiamo di riuscire ad identificare quali sono le lettere di $\Sigma$ presenti in $w$. Dopo la lettura a partire da ogni possibile posizione (stato) del termine $\#$, possiamo identificare un nodo per ciascuna delle situazioni, preparandoci alla lettura dei termini di $w'$:

{
\small
\begin{center}
\begin{tikzpicture}[shorten >=1pt, node distance=2cm, on grid, auto]
    \node[state, initial, initial text=] (init) {$\emptyset$};
    \node[state, right=of init] (b) {$\{b\}$};
    \node[state, above=of b] (a) {$\{a\}$};
    \node[state, below=of b] (c) {$\{c\}$};
    
    \node[state, right=of a] (ab) {$\{a,b\}$};
    \node[state, right=of b] (bc) {$\{b,c\}$};
    \node[state, right=of c] (ac) {$\{a,c\}$};

    \node[state, right=of bc] (abc) {$\{a,b,c\}$};

    % nuovi stati a destra di `abc` (etichette con '#', nomi interni senza '#')
    \node[state, right=of abc, yshift=+4.2cm, fill = blue!5] (ahash) {$\mathrm{a\#}$};
    \node[state, right=of abc, yshift=+3cm, fill = blue!5] (abhash) {$\mathrm{ab\#}$};
    \node[state, right=of abc, yshift=+1.5cm, fill = blue!5] (abchash) {$\mathrm{abc\#}$};
    \node[state, right=of abc, yshift=0cm, fill = blue!5] (bhash) {$\mathrm{b\#}$};
    \node[state, right=of abc, yshift=-1.5cm, fill = blue!5] (bchash) {$\mathrm{bc\#}$};
    \node[state, right=of abc, yshift=-3cm, fill = blue!5] (achash) {$\mathrm{ac\#}$};
    \node[state, right=of abc, yshift=-4.2cm, fill = blue!5] (chash) {$\mathrm{c\#}$};

    \path[->]
    (init) edge node {a} (a)
           edge node {b} (b)
           edge node {c} (c)
    (a) edge [loop left, looseness=3] node {a} (a)
        edge node {b} (ab)
        edge [bend left=3] node {c} (ac)
    (c) edge [loop left, looseness=3] node {c} (c)
        edge node {b} (bc)
        edge [bend right] node {a} (ac)
    (b) edge[loop above, looseness=3] node {b} (b)
        edge [bend right] node {a} (ab)
        edge [bend right] node {c} (bc)
    (ab) edge [loop above, looseness=3] node {a,b} (ab)
         edge node {c} (abc)
    (bc) edge [loop above, looseness=3] node {b,c} (bc)
         edge node {a} (abc)
    (ac) edge[loop above, looseness=3] node {a,c} (ac)
         edge node {b} (abc)
    % archi blu etichettati '#' verso i nuovi stati
    (a)   edge[draw=blue, text=blue] node {$\#$} (ahash)
    (ab)  edge[draw=blue, text=blue] node {$\#$} (abhash)
    (abc) edge[draw=blue, text=blue] node {$\#$} (abchash)
    (b)   edge[draw=blue, text=blue, bend left] node {$\#$} (bhash)
    (ac)  edge[draw=blue, text=blue] node {$\#$} (achash)
    (bc)  edge[draw=blue, text=blue, bend right] node {$\#$} (bchash)
    (c)   edge[draw=blue, text=blue] node {$\#$} (chash);
    ;
\end{tikzpicture}
\end{center}
}


   \vfill
    \noindent\makebox[\linewidth]{\rule{\linewidth}{0.4pt}}
    \newpage
    {\noindent
    \small \textbf{Computabilità, Complessità e Logica - a.a. 2025/26} \hfill \textbf{Approfondimento \EsercitazioneNum} \\[0.2cm]
    \scriptsize Matteo Liotta, \texttt{matteo.liotta@studenti.units.it} \hfill \Data \\[0.2cm]
    \noindent\makebox[\linewidth]{\rule{\linewidth}{0.4pt}} \\[0.3cm]}



Ricordando che per le parole $w$ e $w'$ è sufficiente condividere \textbf{almeno un simbolo} e dal momento che ogni stato azzurro rappresenta intrinsecamente la conoscenza raccolta a seguito alla lettura di $w$, sarà sufficiente leggere una di quelle lettere indicate nello stato per raggiungere l'accettazione -mentre le altre lettere manterranno invariato lo stato. Possiamo concludere la scrittura deterministica dell'automa con la seguente rappresentazione:

{
\small
\begin{center}
\begin{tikzpicture}[shorten >=1pt, node distance=2cm, on grid, auto]
    \node[state, initial, initial text=] (init) {$\emptyset$};
    \node[state, right=of init] (b) {$\{b\}$};
    \node[state, above=of b] (a) {$\{a\}$};
    \node[state, below=of b] (c) {$\{c\}$};
    
    \node[state, right=of a] (ab) {$\{a,b\}$};
    \node[state, right=of b] (bc) {$\{b,c\}$};
    \node[state, right=of c] (ac) {$\{a,c\}$};

    \node[state, right=of bc] (abc) {$\{a,b,c\}$};

    % nuovi stati a destra di `abc` (etichette con '#', nomi interni senza '#')
    \node[state, right=of abc, yshift=+4.6cm, fill = blue!5] (ahash) {$\mathrm{a\#}$};
    \node[state, right=of abc, yshift=+3cm, fill = blue!5] (abhash) {$\mathrm{ab\#}$};
    \node[state, right=of abc, yshift=+1.5cm, fill = blue!5] (abchash) {$\mathrm{abc\#}$};
    \node[state, right=of abc, yshift=-0.1cm, fill = blue!5] (bhash) {$\mathrm{b\#}$};
    \node[state, right=of abc, yshift=-1.5cm, fill = blue!5] (bchash) {$\mathrm{bc\#}$};
    \node[state, right=of abc, yshift=-3cm, fill = blue!5] (achash) {$\mathrm{ac\#}$};
    \node[state, right=of abc, yshift=-4.6cm, fill = blue!5] (chash) {$\mathrm{c\#}$};
    \node[state, right=6of bhash, fill=red!5] (accept) {$accept$};

    \path[->]
    (init) edge node {a} (a)
           edge node {b} (b)
           edge node {c} (c)
    (a) edge [loop left, looseness=3] node {a} (a)
        edge node {b} (ab)
        edge [bend left=3] node {c} (ac)
    (c) edge [loop left, looseness=3] node {c} (c)
        edge node {b} (bc)
        edge [bend right] node {a} (ac)
    (b) edge[loop above, looseness=3] node {b} (b)
        edge [bend right] node {a} (ab)
        edge [bend right] node {c} (bc)
    (ab) edge [loop above, looseness=3] node {a,b} (ab)
         edge node {c} (abc)
    (bc) edge [loop above, looseness=3] node {b,c} (bc)
         edge node {a} (abc)
    (ac) edge[loop above, looseness=3] node {a,c} (ac)
         edge node {b} (abc)
    % archi blu etichettati '#' verso i nuovi stati
    (a)   edge[draw=blue, text=blue, bend left] node {$\#$} (ahash)
    (ab)  edge[draw=blue, text=blue, bend left] node {$\#$} (abhash)
    (abc) edge[draw=blue, text=blue] node {$\#$} (abchash)
    (b)   edge[draw=blue, text=blue, bend left] node {$\#$} (bhash)
    (ac)  edge[draw=blue, text=blue] node {$\#$} (achash)
    (bc)  edge[draw=blue, text=blue, bend right] node {$\#$} (bchash)
    (c)   edge[draw=blue, text=blue, bend right] node {$\#$} (chash)

    (ahash) edge[draw=red, text=red, bend left] node {$a$} (accept)
            edge [loop above, looseness=3] node {b,c} (ahash)
    (abhash) edge[draw=red, text=red, bend left] node {$c$} (accept)
            edge [loop above, looseness=3] node {c} (abhash)
    (abchash) edge[draw=red, text=red] node {$a,b,c$} (accept)

    (bhash) edge[draw=red, text=red] node {$b$} (accept)
            edge [loop above, looseness=3] node {a,c} (bhash)
    (achash) edge[draw=red, text=red, bend right] node {$a,c$} (accept)
            edge [loop above, looseness=3] node {b} (achash)
    (bchash) edge[draw=red, text=red] node {$b$} (accept)
            edge [loop above, looseness=3] node {a} (bchash)
    (chash) edge[draw=red, text=red, bend right] node {$c$} (accept)
            edge [loop above, looseness=3] node {a,b} (chash);

    
    
    ;
\end{tikzpicture}
\end{center}
}

Questo conclude la rappresentazione dell'automa deterministico per il linguaggio indicato.\\[0.3cm]

\underline{\textbf{Caso non deterministico}}\\[0.3cm]
Ricordando che per definizione una parola viene riconosciuta da un automa ND se ammette almeno una computazione accettante, questo permette di considerare una serie di percorsi indipendenti: possiamo \textbf{inizialmente fissare} il termine che ci \textit{aspettiamo} essere in comune tra $w$ e $w'$, attraverso differenti percorsi. Tale scelta sarà effettuata non deterministicamente impiegando $\epsilon$-transizioni: l'automa inizierà con il \textit{guess}, accompagnato da una scelta non deterministica, con lo scopo di dare origine a $n$ percorsi indipendenti semplici, ciascuno volto a verificare l'\textbf{esito} di ogni scommessa.
\\\\
Considerando lo stato iniziale, da $q_0$ vogliamo rendere raggiungibile ciascuno uno dei seguenti:
\begin{itemize}
    \item $S_{x=a}$: stato in cui si scommette che sarà $x=a\in\Sigma$ l'elemento comune tra le parole $w$ e $w'$
    \item $S_{x=b}$: stato in cui si scommette che sarà $x=b\in\Sigma$ l'elemento comune tra le parole $w$ e $w'$
    \item $S_{x=c}$: stato in cui si scommette che sarà $x=c\in\Sigma$ l'elemento comune tra le parole $w$ e $w'$
\end{itemize}

{
\small
\begin{center}
\begin{tikzpicture}[shorten >=1pt, node distance=2cm, on grid, auto]
    \node[state, initial, initial text=] (init)   {$\emptyset$};
    \node[state, right=of init] (Sb) {$S_{x=b}$};
    \node[state, above=1.1cm of b] (Sa) {$S_{x=a}$};
    \node[state, below=1.1cm of b] (Sc) {$S_{x=c}$};

    \path[->]
    (init) edge node {$\epsilon$} (Sa)
    (init) edge node {$\epsilon$} (Sb)
    (init) edge node {$\epsilon$} (Sc);
    
\end{tikzpicture}
\end{center}
}

   \vfill
    \noindent\makebox[\linewidth]{\rule{\linewidth}{0.4pt}}
    \newpage
    {\noindent
    \small \textbf{Computabilità, Complessità e Logica - a.a. 2025/26} \hfill \textbf{Approfondimento \EsercitazioneNum} \\[0.2cm]
    \scriptsize Matteo Liotta, \texttt{matteo.liotta@studenti.units.it} \hfill \Data \\[0.2cm]
    \noindent\makebox[\linewidth]{\rule{\linewidth}{0.4pt}} \\[0.3cm]}


In ciascun ramo raggiungibile (per quanto indipendente), si effettua una scommessa sulla lettera $x$ che ci si aspetta essere l'occorrenza comune in $w$ e $w'$. Due sono le possibilità dunque, raggiunto un ramo:
\begin{itemize}
\item Viene effettivamente letta subito la lettera $x$
\item Non viene letta $x$ e si leggono le altre $x'\in\Sigma_{/x}$, attendendo l'eventuale arrivo della lettura di $x$ su cui si ha scommesso.
\end{itemize}
Qualora la lettura della lettera $x$ avvenisse effettivamente, sarebbe possibile avanzare a uno stato in cui è contenuta implicitamente consapevolezza sull'aver ottenuto il simbolo $x$ nella parola $w$. 

{
\small
\begin{center}
\begin{tikzpicture}[shorten >=1pt, node distance=2cm, on grid, auto]
    \node[state, initial, initial text=] (init)   {$\emptyset$};
    \node[state, right=of init] (Sb) {$S_{x=b}$};
    \node[state, above=2cm of b] (Sa) {$S_{x=a}$};
    \node[state, below=2cm of b] (Sc) {$S_{x=c}$};
    
    \node[state, right=2cm of Sa] (Wa) {$W_{x=a}$};
    \node[state, right=2cm of Sb] (Wb) {$W_{x=b}$};
    \node[state, right=2cm of Sc] (Wc) {$W_{x=c}$};

    \path[->]
    (init) edge node {$\epsilon$} (Sa)
    (init) edge node {$\epsilon$} (Sb)
    (init) edge node {$\epsilon$} (Sc)
    (Sa) edge node {$a$} (Wa)
         edge [loop above, looseness=3] node {$b,c$} (Sa)
    (Sb) edge node {$b$} (Wb)
         edge [loop above, looseness=3] node {$a,c$} (Sb)
    (Sc) edge node {$c$} (Wc)
         edge [loop above, looseness=3] node {$a,b$} (Sc);
    
\end{tikzpicture}
\end{center}
}

A questo punto, è necessaria la lettura del termine $\#$:

{
\small
\begin{center}
\begin{tikzpicture}[shorten >=1pt, node distance=2cm, on grid, auto]
    \node[state, initial, initial text=] (init)   {$\emptyset$};
    \node[state, right=of init] (Sb) {$S_{x=b}$};
    \node[state, above=2cm of b] (Sa) {$S_{x=a}$};
    \node[state, below=2cm of b] (Sc) {$S_{x=c}$};
    
    \node[state, right=2cm of Sa] (Wa) {$W_{x=a}$};
    \node[state, right=2cm of Sb] (Wb) {$W_{x=b}$};
    \node[state, right=2cm of Sc] (Wc) {$W_{x=c}$};

    \node[state, right=2cm of Wa] (Ha) {$W_{\#_{x=a}}$};
    \node[state, right=2cm of Wb] (Hb) {$W_{\#_{x=b}}$};
    \node[state, right=2cm of Wc] (Hc) {$W_{\#_{x=c}}$};
    

    \path[->]
    (init) edge node {$\epsilon$} (Sa)
    (init) edge node {$\epsilon$} (Sb)
    (init) edge node {$\epsilon$} (Sc)
    (Sa) edge node {$a$} (Wa)
         edge [loop above, looseness=3] node {$b,c$} (Sa)
    (Sb) edge node {$b$} (Wb)
         edge [loop above, looseness=3] node {$a,c$} (Sb)
    (Sc) edge node {$c$} (Wc)
         edge [loop above, looseness=3] node {$a,b$} (Sc)
    
    (Wa) edge node {$\#$} (Ha)
         edge [loop above, looseness=3] node {$a,b,c$} (Wa)
    (Wb) edge node {$\#$} (Hb)
         edge [loop above, looseness=3] node {$a,b,c$} (Wb)
    (Wc) edge node {$\#$} (Hc)
         edge [loop above, looseness=3] node {$a,b,c$} (Wc);
    
\end{tikzpicture}
\end{center}
}

Se dunque si ha raggiunto uno stato $W_{\#_{x=...}}$ significa che la il simbolo $x$ è stato trovato nella parola $w$ e si è ora in attesa della lettura della seconda parola $w'$. Se $x$ su cui si ha inizialmente scommesso viene letto in $w'$, allora sarà possibile continuare verso lo stato finale:

{
\small
\begin{center}
\begin{tikzpicture}[shorten >=1pt, node distance=2cm, on grid, auto]
    \node[state, initial, initial text=] (init)   {$\emptyset$};
    \node[state, right=of init] (Sb) {$S_{x=b}$};
    \node[state, above=2cm of b] (Sa) {$S_{x=a}$};
    \node[state, below=2cm of b] (Sc) {$S_{x=c}$};
    
    \node[state, right=2cm of Sa] (Wa) {$W_{x=a}$};
    \node[state, right=2cm of Sb] (Wb) {$W_{x=b}$};
    \node[state, right=2cm of Sc] (Wc) {$W_{x=c}$};

    \node[state, right=2cm of Wa] (Ha) {$W_{\#_{x=a}}$};
    \node[state, right=2cm of Wb] (Hb) {$W_{\#_{x=b}}$};
    \node[state, right=2cm of Wc] (Hc) {$W_{\#_{x=c}}$};

    \node[state, right=2cm of Ha, accepting] (OKa) {${x=a}$};
    \node[state, right=2cm of Hb, accepting] (OKb) {${x=b}$};
    \node[state, right=2cm of Hc, accepting] (OKc) {${x=c}$};
    

    \path[->]
    (init) edge node {$\epsilon$} (Sa)
    (init) edge node {$\epsilon$} (Sb)
    (init) edge node {$\epsilon$} (Sc)
    (Sa) edge node {$a$} (Wa)
         edge [loop above, looseness=3] node {$b,c$} (Sa)
    (Sb) edge node {$b$} (Wb)
         edge [loop above, looseness=3] node {$a,c$} (Sb)
    (Sc) edge node {$c$} (Wc)
         edge [loop above, looseness=3] node {$a,b$} (Sc)
    
    (Wa) edge node {$\#$} (Ha)
         edge [loop above, looseness=3] node {$a,b,c$} (Wa)
    (Wb) edge node {$\#$} (Hb)
         edge [loop above, looseness=3] node {$a,b,c$} (Wb)
    (Wc) edge node {$\#$} (Hc)
         edge [loop above, looseness=3] node {$a,b,c$} (Wc)

    (Ha) edge node {$a$} (OKa)
         edge [loop above, looseness=3] node {$a,b,c$} (Ha)
    (Hb) edge node {$b$} (OKb)
         edge [loop above, looseness=3] node {$a,b,c$} (Hb)
    (Hc) edge node {$c$} (OKc)
         edge [loop above, looseness=3] node {$a,b,c$} (Hc);
    
\end{tikzpicture}
\end{center}
}

   \vfill
    \noindent\makebox[\linewidth]{\rule{\linewidth}{0.4pt}}
    \newpage
    {\noindent
    \small \textbf{Computabilità, Complessità e Logica - a.a. 2025/26} \hfill \textbf{Approfondimento \EsercitazioneNum} \\[0.2cm]
    \scriptsize Matteo Liotta, \texttt{matteo.liotta@studenti.units.it} \hfill \Data \\[0.2cm]
    \noindent\makebox[\linewidth]{\rule{\linewidth}{0.4pt}} \\[0.3cm]}
 
Questo conclude dunque la realizzazione dell'automa non deterministico che riconosca il linguaggio $L$. 
\\\\

L'intuizione vista a lezione alla base della "\textit{scommessa}" non deterministica si può ora astrarre con più facilità dall'esempio: poiché in un automa non deterministico si definisce accettata una parola tale da generare una computazione accettante su almeno uno dei "rami", possiamo considerare $n$ rami separati, paralleli, in cui sicuramente almeno uno porterà ad accettazione se verificate le condizioni del linguaggio. \\ 
Per quanto abituati a considerare un singolo valore $x$ per volta, in realtà nel momento in cui stiamo considerando la computazione su


{
\small
\begin{center}
\begin{tikzpicture}[shorten >=1pt, node distance=2cm, on grid, auto]
    \node[state, initial, initial text=] (init)   {$\emptyset$};
    \node[state, right=of init] (Sb) {$S_{x=b}$};
    \node[state, above=2cm of b] (Sa) {$S_{x=a}$};
    \node[state, below=2cm of b] (Sc) {$S_{x=c}$};
    
    \node[state, right=2cm of Sa] (Wa) {$W_{x=a}$};
    \node[state, right=2cm of Sb] (Wb) {$W_{x=b}$};
    \node[state, right=2cm of Sc] (Wc) {$W_{x=c}$};

    \node[state, right=2cm of Wa] (Ha) {$W_{\#_{x=a}}$};
    \node[state, right=2cm of Wb] (Hb) {$W_{\#_{x=b}}$};
    \node[state, right=2cm of Wc] (Hc) {$W_{\#_{x=c}}$};

    \node[state, right=2cm of Ha, accepting] (OKa) {${x=a}$};
    \node[state, right=2cm of Hb, accepting] (OKb) {${x=b}$};
    \node[state, right=2cm of Hc, accepting] (OKc) {${x=c}$};
    

    \path[->]
    (init) edge node {$\epsilon$} (Sa)
    (init) edge node {$\epsilon$} (Sb)
    (init) edge node {$\epsilon$} (Sc)
    (Sa) edge node {$a$} (Wa)
         edge [loop above, looseness=3] node {$b,c$} (Sa)
    (Sb) edge node {$b$} (Wb)
         edge [loop above, looseness=3] node {$a,c$} (Sb)
    (Sc) edge node {$c$} (Wc)
         edge [loop above, looseness=3] node {$a,b$} (Sc)
    
    (Wa) edge node {$\#$} (Ha)
         edge [loop above, looseness=3] node {$a,b,c$} (Wa)
    (Wb) edge node {$\#$} (Hb)
         edge [loop above, looseness=3] node {$a,b,c$} (Wb)
    (Wc) edge node {$\#$} (Hc)
         edge [loop above, looseness=3] node {$a,b,c$} (Wc)

    (Ha) edge node {$a$} (OKa)
         edge [loop above, looseness=3] node {$a,b,c$} (Ha)
    (Hb) edge node {$b$} (OKb)
         edge [loop above, looseness=3] node {$a,b,c$} (Hb)
    (Hc) edge node {$c$} (OKc)
         edge [loop above, looseness=3] node {$a,b,c$} (Hc);
    
\end{tikzpicture}
\end{center}
}

tutti ed $n$ i processi coesistono paralleli. \\\\Infine, poiché esistono $n$ rami paralleli, ognuno specializzato per una singola lettera, siamo certi che una parola in cui $w$ e $w'$ condividono $x\in\Sigma$ saranno accettate.
%nel momento in cui ci interrogheremo sulla possibilità di riconoscere quella parola da parte di un automa ND, dobbiamo considerare il comportamento su ciascuno dei percorsi che esso ammette; se presente uno tale da portare alla accettazione, allora la parola sarà nel linguaggio di interesse.

La differenza e il vantaggio risiedono dunque nella possibilità di considerare $n$ possibilità e scommesse (raggiungibili non-deterministicamente) che verifichino indipendentemente ciascuna possibilità, attraverso strutture più semplici.\\

Risulta allora evidente come la scrittura non deterministica sia più compatta rispetto alla precedente, in termini di cardinalità dell'insieme degli stati richiesti. 


\vfill
\noindent\makebox[\linewidth]{\rule{\linewidth}{0.1pt}}

\end{document}
